\documentclass{article}
\usepackage{/globals/math}
\begin{document}
\section{Zahlenfolgen}
\subsection{Summen}
Mit dem Summenzeichen $\Sigma$ kann einfach notiert werden, dass eine bestimmte Anzahl an Summanden einer bestimmten Reihe addiert werden sollen.
\[
 \sum_{k=m}^n a_k = a_m + a_{m+1} + \dots + a_n
\]

\subsection{Produkte}
Das gleiche, nur mit Multiplikation anstelle von Addition, kann mit dem Produktzeichen $\Pi$ passieren.
\[
 \prod_{k=m}^n a_k = a_m * a_{m+1} * \dots * a_n
\]

\subsection{Anwendungen}
Die Fakultät $n!$ beschreibt wie viele Möglichkeiten es gibt $n$ Objekte zu sortieren. Hier gilt
\[
 n! = \prod_{k=1}^n k
\]

Wird betrachtet wie viele Möglichkeiten es gibt $k$ Elemente von $n$ Elementen auszuwählen, wobei nur betrachtet wird, welche Elemente ausgewählt wurden, nicht in welcher Reihenfolge, so wird der Binomialkoeffizient genutzt.
\[
 \binom{n}{k} = \frac{n!}{k! * (n-k)!}
\] 

\end{document}